\section{Introduction}
\label{sec:intro}
3D reconstruction stands as a cornerstone technology for augmented/virtual reality, autonomous systems, and digital twin applications. Traditional pipelines~\cite{schonberger2016structure, mur2015orb} relying on Structure-from-Motion (SfM) and Multi-View Stereo (MVS) achieve remarkable precision through iterative optimization, but still suffer from computational complexity and fragile performance in textureless regions. The recent paradigm shift toward feed-forward approaches~\cite{dust3r_cvpr24, leroy2024mast3r, wang2025vggt, keetha2025mapanything, wang2025pi3} has revolutionized this landscape, enabling real-time, generalizable 3D reconstruction by training forward-pass deep networks with a large amount of data. The powerful baselines like MASt3R~\cite{leroy2024mast3r} and VGGT~\cite{wang2025vggt} have demonstrated unprecedented capability in joint camera poses and dense geometry prediction, establishing a strong foundation with robust geometry prior for downstream tasks, especially novel view synthesis (NVS)~\cite{wang2025volsplat, jiang2025anysplat, ye2024nopo} with 3D Gaussian Splatting (3DGS)~\cite{kerbl3Dgaussians}.

Despite recent advances, feed-forward 3DGS~\cite{wang2025volsplat, jiang2025anysplat, ye2024nopo} remains fundamentally limited for extrapolated view renderings when input views are sparse. Specifically, when rendering extrapolated views has large area outside the input boundaries, the rendered results from 3DGS exhibits severe artifacts such as incomplete geometry, texture corruption, and large blank areas where no observations exist as in \cref{fig:teaser}. Previous latent diffusion models~\cite{Rombach_2022_SD2, suvorov2021resolutionLaMa} excel at 2D inpainting, they lack knowledge of the underlying 3D structure, producing {multi-view inconsistent} hallucinations that violate scene geometry when projected across viewpoints. While, existing works~\cite{wu2025difix3d+,zhong2025taming,fischer2025flowr,wei2025gsfix3d,yin2025gsfixer} propose some novel techniques to address the rendering artifacts of 3DGS, they predominantly employ diffusion models as post-processing tools for 2D rendering correction. Although they improves perceptual quality within observed regions, the generation results still lacks geometric awareness, resulting in catastrophic failures such as generating inconsistent depth, repeating textures, blurring artifacts when encountering large extrapolated areas.


Stemming from the fundamental ill-posedness of view extrapolation, we identify that the feedback loop between generation and reconstruction remains underexplored. Prior work treats diffusion refinement as a 2D post-process, overlooking a critical limitation that latent diffusion models lack 3D geometric priors knowledge. To bridge this gap, we introduce Leveling3D, the first framework that tightly couples feed-forward 3D reconstruction with geometry-conditioned diffusion generation. Built on AnySplat~\cite{jiang2025anysplat} and Stable Diffusion 2~\cite{Rombach_2022_SD2}, Leveling3D bridges reconstruction and generation through a geometry-aware adapter (Fig.~\ref{fig:overview}). Our key insights are: (1) transfering 3D geometric priors (tokens aggregated from multi-view inputs by VGGT~\cite{wang2025vggt}) to condition the diffusion model, grounding 2D generation in 3D structure; (2) the geometrically consistent extrapolated views generated by diffusion model are fed back to refine the 3DGS representation. This reconstruction-generation-refinement produces photorealistic details in underconstrained regions while preserving multi-view geometric consistency.
Our proposed pipeline comprises three key components: (i) a lightweight \emph{geometry-aware leveling adapter} that injects 3D geometric priors into diffusion via cross-attention; (ii) a \emph{palette filtering strategy} that ensure a well-behaved, approximately normal training distribution that prevents mode collapse (iii) a \emph{test-time masking refinement} that ensures seamless boundaries between reconstructed and generated regions. We achieve state-of-the-art performance on public unseen datasets across multiple scenes, including both the NVS rendering quality and geometry consistency, demonstrating that our geometry-aware approach significantly outperforms both vanilla diffusion model and other 3DGS fixing methods based on diffusion model.
Our contributions can be summarized as:
\begin{itemize}
    \item Leveling3D, a novel framework to tightly couple feed-forward 3D reconstruction with generation.
    \item Leveling adapter with a palette filtering training strategy and mask refinement that enables diffusion control with geometric priors for complex 3D-aware extrapolated view refinement.
    % \item A test-time masking refinement that produces artifact-free boundaries between reconstructed and generated regions.
    \item A refinement pipeline that leverages enhanced extrapolated views feedback to improve 3D representation quality.
    \item Extensive experiments demonstrate that Leveling3D outperforms existing SOTA on public datasets on both novel-view synthesis and depth estimation.
\end{itemize}



